%%%%%%%%%%  scan idx 198  |  folio 185  %%%%%%%%%%
liser au voisinage de $s$, que les $M_i$ (resp. $p_i$) sont les fibres en $s$ de sous-$\mathcal O_S$-Modules cohérents $F_i$ (resp. Idéaux cohérents $\underline{p}_i$) de $F$ (resp. $\mathcal O_S$) tels que l'on ait
\[
 F=F_0\supset F_1\supset\ldots\supset F_q=0
\]
et $F_i/F_{i+1}\simeq\mathcal O_S/\underline{p}_i$. Prouver que la suite $(E_n)$ est stationnaire au voisinage de $s$ équivaut évidemment à prouver que les suites induites sur chaque cran de la filtration $(F_i)$ de $F$ le sont. De sorte qu'on peut supposer que $F=\mathcal O_S$, et que $\mathcal O_{S,s}$ est un anneau intègre, les $E_n$ formant une suite croissante d'Idéaux de $\mathcal O_S$ tels que $(E_n)_s=0$. Comme $\mathcal O_s$ est intègre, donc en particulier réduit, on peut en outre supposer grâce au théorème d'Oka ([10] IV \S3), quitte à se localiser en $s$, que $S$ est \emph{réduit}. Montrons qu'il existe un voisinage ouvert $U$ de $s$ tel que $E_n|U=0$ pour $n$ assez grand. Considérons pour cela le normalisé $\widetilde S$ de $S$ ([10] IV \S4). On sait (loc. cit.) que $S$ est le spectre analytique ([10] II \S3) d'une $\mathcal O_S$-Algèbre finie $\widetilde{\mathcal O}_S$ dont la fibre en chaque $t\in S$ est la clôture intégrale de $\mathcal O_{S,t}$ (cf. note $(^2)$ de (2.4.8)). On a $E_n\subset E_n\cdot\widetilde{\mathcal O}_S$, donc il suffit de prouver que $E_n\cdot\widetilde{\mathcal O}_S=0$ pour $n$ assez grand. Or on a (10 II \S3 Prop.8)
\[
 E_n\cdot\widetilde{\mathcal O}_S = p_*(\widetilde{E_n\cdot\widetilde{\mathcal O}_S}),
\]
où $p:\widetilde S\longrightarrow S$ est la projection canonique, et $\widetilde{E_n\cdot\widetilde{\mathcal O}_S}$ l'Idéal de $\mathcal O_{\widetilde S}$ défini canoniquement par $E_n\cdot\widetilde{\mathcal O}_S$. Il suffit donc de montrer qu'il existe un voisinage ouvert $V$ de $s$ dans $\widetilde S$ tel que $\widetilde{E_n\cdot\widetilde{\mathcal O}_S}|U=0$ pour $n$ assez grand, ce qui nous ramène au cas où $S$ est \emph{normal}, ce que nous supposerons maintenant. Comme $(E_n)_s=0$, l'immersion fermée $i_n:S_n\longrightarrow S$ définie par $E_n$ est \emph{étale} en $s$. Mais, $S$ étant normal, l'ensemble des points où $i_n$ est étale est à la fois ouvert et fermé (voir lemme ci-dessous), donc son image dans $S$ est une partie de $S$ à la fois ouverte et fermée. Donc la restriction de $E_n$ à la composante connexe de $s$ dans $S$ est nulle ; cette composante connexe étant un ensemble ouvert, puisque $S$ est localement connexe, on a gagné. Il reste donc à établir le

